\section{The Generalized Continuum Hypothesis in $L$}

The last axiom on Gödel's list takes the most work, and its proof uses an idea particular to $L$.

\begin{boxtheorem}[Working in $L$] \label{Ch4:Thm:GCH_In_L}
    $(\GCH)^L$. Ie,
    \begin{align*}
        \parenth{\text{For all infinite cardinals $\kappa$, } \kappa^+ = 2^{\kappa}}^L
    \end{align*}
\end{boxtheorem}
\begin{proof}[Proof, working in $L$ and keeping in mind that $L^L = L$]
    Let $\lambda$ be an infinite cardinal. We will show that $H_{\lambda} = L_{\lambda}$. This will be enough because if we have any infinite cardinal $\kappa$, setting $\lambda = \kappa^+$, we get that
    \begin{align*}
        \Powset(\kappa) \ssq H_{\lambda} = L_{\lambda}
    \end{align*}
    meaning that
    \begin{align*}
        2^{\kappa} = \abs{\Powset(\kappa)} \leq \abs{L_{\lambda}} = \abs{\lambda} = \lambda = \kappa^+
    \end{align*}
    % [CORRECTED] was: ``always true in $V$, in $L$, everywhere'' --- $\kappa^+ \leq 2^{\kappa}$ needs choice (ZF alone does not prove $\aleph_1 \leq 2^{\aleph_0}$), and $V$ is only assumed to satisfy $\ZF$ here; it does hold in $L$, which is all the proof needs
    And of course the other inequality is always true in $L$, as in any model of $\ZFC$.

    So now let's prove that $H_{\lambda} = L_{\lambda}$ for all infinite cardinals $\lambda$.

    Fix an infinite cardinal $\lambda$. We already know that
    \begin{align*}
        H_{\omega} = V_{\omega} = L_{\omega}
    \end{align*}
    so we may assume that $\lambda$ is uncountable.
    \begin{itemize}
        % [CORRECTED] was: ``Pick $\alpha < \lambda$ such that $z \in L_{\alpha}$'' --- $\abs{L_{\alpha}} = \abs{\alpha}$ fails for finite $\alpha \geq 3$ (eg $\abs{L_3} = 4$); as $\lambda > \omega$ we may take $\alpha \geq \omega$
        \item \underline{$L_{\lambda} \ssq H_{\lambda}$.} Let $z \in L_{\lambda}$. Pick $\alpha$ with $\omega \leq \alpha < \lambda$ such that $z \in L_{\alpha}$. Then certainly $\trcl{\set{z}} \ssq L_{\alpha}$. So
        \begin{align*}
            \abs{\trcl{\set{z}}} \leq \abs{L_{\alpha}} = \abs{\alpha}\leq \alpha < \lambda
        \end{align*}
        So $z \in H_{\lambda}$.

        \item \underline{$H_{\lambda} \ssq L_{\lambda}$.} Let $z \in H_{\lambda}$. Recall our finite subtheory $T$ of $\ZFwoP$, for whose transitive models $\alpha \mapsto L_{\alpha}$ and $\alpha \mapsto <_{\alpha}$ are absolute, maybe augmented with $V = L$.
        By the reflection theorem scheme, there's some $\beta$ such that $z \in L_{\beta}$ and $\sigma^{L_{\beta}}$ for all $\sigma \in T$ and $(V = L)^{L_{\beta}}$. So $L_{\beta} \models T$ and $(V = L)^{L_{\beta}}$.

        Observe: $\trcl{\set{z}} \ssq L_{\beta}$ and $\abs{\trcl{\set{z}}} < \lambda$. By Downwards Löwenheim-Skolem, we have some $X \preceq L_{\beta}$ with $\abs{X} < \lambda$ and $\trcl{\set{z}} \ssq X$. Let $G : X \simeq M$ be the Mostowski Collapse of $X$, so $M$ is a transitive model of $T$. Since $\trcl{\set{z}} \ssq X$, we know that
        \begin{align*}
            G \restriction \trcl{\set{z}} = \id \restriction \trcl{\set{z}}
        \end{align*}
        In particular, $G(z) = z \in M$.

        As $L_{\beta} \models V = L$, we also have $M \models V = L$. In other words, $(V = L)^M$. This says that $V^M = L^M$, which says $M = L_{\alpha}$ where $\alpha = \OR \cap M$. Therefore, $z \in L_{\alpha}$ and $\abs{X} = \abs{M} = \abs{L_{\alpha}} = \abs{\alpha}$, and this is less than $\lambda$, so $\alpha < \lambda$.
    \end{itemize}
\end{proof}

% [CORRECTED] was: ``some nice level of the $L$ hierarchy (which we get by downwards Löwenheim-Skolem), we take an elementary substructure'' --- in the proof the level comes from the reflection theorem scheme and the elementary substructure from Löwenheim-Skolem
We do something very interesting in the above proof: we take some nice level of the $L$ hierarchy (which we get by the reflection theorem scheme), we take an elementary substructure (which we get by downwards Löwenheim-Skolem), then take its Mostowski collapse, and we end up with \textit{another} level of the $L$ hierarchy.

It turns out that this is not true of the $V$- or the $H$-hierarchy. The phenomenon has a name.

\begin{boxdefinition}[Condensation]
    The phenomenon that the Mostowski collapse of an elementary substructure of a suitable level of the $L$-hierarchy---such as the $L_{\beta}$ in the proof above, with $L_{\beta} \models T$ and $(V = L)^{L_{\beta}}$---is again a level of the $L$-hierarchy is known as \textbf{condensation}.
\end{boxdefinition}

The condensation used in that proof can be stated once and for all, for any level of the $L$-hierarchy satisfying enough of $\ZFwoP$.

\begin{boxlemma}[Condensation Lemma]
    Let $T$ be some finite subtheory of $\ZFwoP$ such that for transitive class models of $T$,
    \begin{align*}
        A \mapsto \D(A)
        \qquad \qquad
        \alpha \mapsto L_{\alpha}
        \qquad \qquad
        \alpha \mapsto <_{\alpha}
        \qquad \qquad
        L
        \qquad \qquad
        <_L
    \end{align*}
    are all absolute. By the reflection theorem (scheme), $\setst{\beta}{L_{\beta} \models T}$ is unbounded in $\OR$. Given $L_{\beta} \models T$ and $X \preceq L_{\beta}$, if $M \simeq X$ and $M$ is transitive, then $M = L_{\alpha}$, where $\alpha = \OR \cap M$.
\end{boxlemma}

Here's another way to view the condensation lemma.

Given a transitive $M$ and an elementary embedding $\pi : M \inj L_{\beta}$ (for some $\beta$ as above), we conclude that there is some $\alpha \leq \beta$ such that $M = L_{\alpha}$. In the statement above, $X = \ran(\pi)$, with $\pi$ being the inverse of the Mostowski collapse of $X$.
